\documentclass[11pt]{article}

%------------------------------------------------
% Packages
%------------------------------------------------
\usepackage[a4paper,margin=2.5cm]{geometry}
\usepackage{amsmath,amssymb,amsthm}
\usepackage{mathtools}
\usepackage{enumitem}
\usepackage{hyperref}

\hypersetup{
    colorlinks=true,
    linkcolor=blue,
    citecolor=blue,
    urlcolor=blue
}

%------------------------------------------------
% Title Information
%------------------------------------------------

\title{\bfseries Locally-averaged McCormick relaxations for discretization-regularized inverse problems
}

\author{
Barbara Kaltenbacher\thanks{Department of Mathematics, University of Klagenfurt,
\texttt{barbara.kaltenbacher@aau.at}}
\and
Paul Manns\thanks{
Department of Mathematics, TU Dortmund University, 
\texttt{paul.manns@tu-dortmund.de}}
}

\date{}

%------------------------------------------------

\begin{document}

\maketitle

\vspace{-1em}

\noindent
\textbf{Keywords:}
inverse problems, regularization by discretization, McCormick relaxation

\vspace{0.5em}


\vspace{1em}

%%%%%%%%%%%%%%%%%%%%%%%%%%%%%%%%%%%%%%%%%%%%%%%%%%
\section*{Abstract}

As a prototypical coefficient inverse problem, we consider identification of the potential in an elliptic PDE from indirect observations and its regularization by discretization in pre-image space -- the so-called least squares version of regularization by projection. We are particularly interested in a situation where the coefficient only takes values in a discrete and finite set. In doing so, we seek to minimize the data misfit over the finite dimensional subspace of piecewise constant function on a given partition of the domain, under the PDE as well as the discrete valuedness as constraints.
The resulting algorithmic challenge is twofold:
* the discretization level must be adapted to the noise level to allow us to stably recover coefficient and state in this nonlinear inverse problem;
* the discretization-regularized minimization problem must be solved to global optimality.
For the former, we establish a general result on the least squares method for nonlinear ill-posed operator equations.
The latter is not possible by gradient-based nonlinear programming algorithms due to non-convexity.
We therefore need a global optimization algorithm like (spatial) branch-\&-bound, which requires means to compute (approximate) primal upper and dual lower bounds. 
In order to obtain dual bounds, we employ a McCormick relaxation, combined with local averaging \cite{LeyfferManns2025mccormick}, which reduces the number of degrees of freedom and makes the problem amenable to optimization based bound tightening.


%%%%%%%%%%%%%%%%%%%%%%%%%%%%%%%%%%%%%%%%%%%%%%%%%%
\vspace{1em}

\noindent
\textbf{Acknowledgements.}
This work was supported by the Austrian Science Fund (FWF) under grant
[10.55776/F100800]

%%%%%%%%%%%%%%%%%%%%%%%%%%%%%%%%%%%%%%%%%%%%%%%%%%
\vspace{1em}

\begin{thebibliography}{9}

\bibitem{IP_McCormick}
Barbara Kaltenbacher and Paul Manns.
\newblock {Locally-averaged McCormick relaxations for
  discretization-regularized inverse problems }.
\newblock {\em Inverse Problems}, 2026.
\newblock see also arXiv:2601.01995 [math.NA].

\bibitem{LeyfferManns2025mccormick}
Sven Leyffer and Paul Manns.
\newblock {McCormick} envelopes in mixed-integer {PDE}-constrained
  optimization.
\newblock {\em Mathematical Programming}, pages 1--46, 2025.

\end{thebibliography}

\end{document}
